\documentclass[11pt]{article}

\usepackage[margin=0.75in]{geometry}
\usepackage{amsmath,amssymb}
\usepackage{booktabs}
\usepackage{array}
\usepackage{enumitem}
\usepackage[hidelinks]{hyperref}

\setlist[itemize]{leftmargin=1.4em}
\setlist[enumerate]{leftmargin=1.6em}

\title{SYSC 5108 Exam Study\\Assignment 2, Question 7}
\author{Jesse Levine}
\date{}

\begin{document}
\maketitle

\section*{Question Summary}

Question 7 uses histogram features for small \(8\times 8\) handwritten images.
Each image is represented by counts of gray levels \(0,\ldots,7\), denoted
\(\text{GL-0},\ldots,\text{GL-7}\). A logistic regression model classifies digits
as either \(0\) or \(1\). Some predictions, errors, squared errors, and
least-squared-error weight-update terms are missing. Then the question asks for
updated weights using LSE, predictions for digits 7 and 8, and a one-sample MLE
update.

\section*{Course and Textbook Cross-References}

\begin{itemize}
    \item \textbf{Chapter 3 slides, Logistic Regression Model:} for \(K=2\),
    logistic regression uses a linear score \(z=\mathbf{w}^T\mathbf{x}\) and
    probability \(p=\sigma(z)\).
    \item \textbf{Chapter 3 slides, Logistic Sigmoid Function:}
    \[
        \sigma(z)=\frac{1}{1+\exp(-z)}.
    \]
    \item \textbf{Chapter 3 slides, Alternative Notation:} the MLE or
    cross-entropy gradient has the clean form
    \[
        -\nabla_w J_{\text{MLE}} = (y-\sigma(z))x.
    \]
    \item \textbf{Question 6 result:} the LSE-on-probabilities update direction is
    \[
        -\nabla_w J_{\text{LSE}} = (y-\sigma(z))\sigma(z)(1-\sigma(z))x.
    \]
    \item \textbf{Goodfellow, Bengio, and Courville, \emph{Deep Learning}:}
    sigmoid outputs model Bernoulli probabilities; cross-entropy/MLE is generally
    preferred over MSE/LSE for sigmoid outputs because LSE adds an extra sigmoid
    derivative factor that can saturate.
\end{itemize}

\section*{Given Model}

Use an intercept feature \(x[0]=1\). For one image,
\[
    z = w[0] + \sum_{j=1}^{8} w[j]x[j],
    \qquad
    \sigma(z)=\frac{1}{1+\exp(-z)}.
\]

The original weights are:

\begin{center}
\begin{tabular}{rrrrrrrrr}
\toprule
\(w[0]\) & \(w[1]\) & \(w[2]\) & \(w[3]\) & \(w[4]\) & \(w[5]\) & \(w[6]\) & \(w[7]\) & \(w[8]\) \\
\midrule
0.309 & 0.100 & -0.152 & -0.163 & 0.191 & -0.631 & -0.716 & -0.478 & -0.171 \\
\bottomrule
\end{tabular}
\end{center}

\section*{Training Histograms}

\begin{center}
\scriptsize
\begin{tabular}{rrrrrrrrrr}
\toprule
ID & GL-0 & GL-1 & GL-2 & GL-3 & GL-4 & GL-5 & GL-6 & GL-7 & Digit \\
\midrule
0 & 31 & 3 & 6 & 2 & 7 & 5 & 6 & 4 & 0 \\
1 & 37 & 3 & 1 & 4 & 1 & 3 & 2 & 13 & 1 \\
2 & 31 & 3 & 4 & 1 & 8 & 7 & 3 & 7 & 0 \\
3 & 38 & 2 & 3 & 0 & 1 & 1 & 5 & 14 & 1 \\
4 & 31 & 5 & 3 & 2 & 5 & 2 & 5 & 11 & 0 \\
5 & 32 & 6 & 3 & 2 & 1 & 1 & 5 & 14 & 1 \\
6 & 31 & 3 & 5 & 2 & 3 & 6 & 2 & 12 & 0 \\
7 & 31 & 4 & 3 & 4 & 1 & 5 & 5 & 11 & 0 \\
8 & 38 & 4 & 2 & 2 & 2 & 4 & 4 & 8 & 1 \\
9 & 38 & 3 & 2 & 3 & 4 & 4 & 1 & 9 & 1 \\
\bottomrule
\end{tabular}
\end{center}

\section*{Part (a): Missing Model Predictions}

For each row, compute
\[
    z_i = w^T x_i,
    \qquad
    \sigma(z_i)=\frac{1}{1+\exp(-z_i)}.
\]

The missing predictions are:

\begin{center}
\begin{tabular}{rrrr}
\toprule
ID & \(z_i\) & \(\sigma(z_i)\) & Predicted class \\
\midrule
1 & -1.8040 & 0.1414 & 0 \\
4 & -6.3160 & 0.0018 & 0 \\
6 & -6.6770 & 0.0013 & 0 \\
\bottomrule
\end{tabular}
\end{center}

\section*{Part (b): Missing Error and Squared Error Values}

The table uses
\[
    \text{error}_i = y_i-\sigma(z_i),
    \qquad
    \text{s.e.}_i = \left(y_i-\sigma(z_i)\right)^2.
\]

\begin{center}
\begin{tabular}{rrr}
\toprule
ID & Missing value & Value \\
\midrule
0 & s.e. & \(0.0000\) \\
5 & error & \(0.9744\) \\
7 & error & \(-0.0045\) \\
8 & s.e. & \(0.9587\) \\
\bottomrule
\end{tabular}
\end{center}

\textbf{Check:} the \(0.9744\) error belongs to ID 5, because
\[
    1-0.0256=0.9744.
\]

\section*{Part (c): Missing \(w[j]\) Error Values}

The table defines the LSE update direction as
\[
    w[j]\text{ error}
    =
    -\frac{\partial L}{\partial w[j]}
    =
    (y_i-\sigma(z_i))\sigma(z_i)(1-\sigma(z_i))x_i[j].
\]

The missing \(w[j]\) error values are:

\begin{center}
\begin{tabular}{rrr}
\toprule
ID & Weight index & Missing value \\
\midrule
3 & \(w[0]\) & \(0.0503\) \\
7 & \(w[2]\) & \(-0.0001\) \\
2 & \(w[8]\) & \(0.0000\) \\
\bottomrule
\end{tabular}
\end{center}

\section*{Part (d): New Weights Using LSE}

Use learning rate \(\eta=0.01\) and a batch update over all ten training samples:
\[
    w_{\text{new}}[j]
    =
    w_{\text{old}}[j]
    +
    \eta\sum_{i=0}^{9}
    (y_i-\sigma(z_i))\sigma(z_i)(1-\sigma(z_i))x_i[j].
\]

The summed update directions and new weights are:

\begin{center}
\begin{tabular}{rrr}
\toprule
Weight & \(\sum_i w[j]\text{ error}_i\) & New weight \\
\midrule
\(w[0]\) & 0.2263 & 0.3113 \\
\(w[1]\) & 8.3484 & 0.1835 \\
\(w[2]\) & 0.7212 & -0.1448 \\
\(w[3]\) & 0.4229 & -0.1588 \\
\(w[4]\) & 0.5878 & 0.1969 \\
\(w[5]\) & 0.3287 & -0.6277 \\
\(w[6]\) & 0.5771 & -0.7102 \\
\(w[7]\) & 0.6887 & -0.4711 \\
\(w[8]\) & 2.8061 & -0.1429 \\
\bottomrule
\end{tabular}
\end{center}

\section*{Part (e): Predictions for Digits 7 and 8 Using LSE Weights}

The new query histograms are:

\begin{center}
\begin{tabular}{rrrrrrrrr}
\toprule
Digit & GL-0 & GL-1 & GL-2 & GL-3 & GL-4 & GL-5 & GL-6 & GL-7 \\
\midrule
7 & 35 & 1 & 5 & 4 & 5 & 2 & 4 & 8 \\
8 & 30 & 6 & 2 & 0 & 5 & 4 & 4 & 13 \\
\bottomrule
\end{tabular}
\end{center}

Using the LSE-updated weights from part (d):

\begin{center}
\begin{tabular}{rrrr}
\toprule
True digit & \(z\) & \(\sigma(z)\) & Predicted class \\
\midrule
7 & -1.0049 & 0.2680 & 0 \\
8 & -5.0926 & 0.0061 & 0 \\
\bottomrule
\end{tabular}
\end{center}

\textbf{Comment:} both out-of-distribution examples are predicted as class \(0\).
This is not very meaningful, because the model was trained only to distinguish
digits \(0\) and \(1\). Digits \(7\) and \(8\) are outside the training classes, so the
outputs should only be interpreted as ``more like 0 than 1'' according to this
particular histogram-based model.

\section*{Part (f): MLE Update Using Only Sample ID 0}

For MLE or cross-entropy, use the update direction
\[
    -\nabla_w J_{\text{MLE},i}
    =
    (y_i-\sigma(z_i))x_i.
\]
The question asks to use only sample ID 0. For ID 0,
\[
    y_0=0,
    \qquad
    \sigma(z_0)=0.00010184,
\]
so
\[
    y_0-\sigma(z_0) = -0.00010184.
\]
With \(\eta=0.01\),
\[
    w_{\text{new}} = w_{\text{old}} + 0.01(y_0-\sigma(z_0))x_0.
\]

\begin{center}
\begin{tabular}{rr}
\toprule
Weight & MLE one-sample updated value \\
\midrule
\(w[0]\) & 0.3090 \\
\(w[1]\) & 0.1000 \\
\(w[2]\) & -0.1520 \\
\(w[3]\) & -0.1630 \\
\(w[4]\) & 0.1910 \\
\(w[5]\) & -0.6310 \\
\(w[6]\) & -0.7160 \\
\(w[7]\) & -0.4780 \\
\(w[8]\) & -0.1710 \\
\bottomrule
\end{tabular}
\end{center}

The weights appear unchanged to four decimal places because ID 0 is already
classified almost perfectly as \(0\), so \(y_0-\sigma(z_0)\) is extremely small.

Repeating part (e) with these one-sample MLE-updated weights:

\begin{center}
\begin{tabular}{rrrr}
\toprule
True digit & \(z\) & \(\sigma(z)\) & Predicted class \\
\midrule
7 & -4.2623 & 0.0139 & 0 \\
8 & -8.0841 & 0.0003 & 0 \\
\bottomrule
\end{tabular}
\end{center}

\section*{Final Answer}

\begin{itemize}
    \item Missing predictions:
    \[
        \sigma(z_1)=0.1414,\qquad
        \sigma(z_4)=0.0018,\qquad
        \sigma(z_6)=0.0013.
    \]
    \item Missing error/squared-error values:
    \[
        \text{s.e.}_0=0.0000,\quad
        \text{error}_5=0.9744,\quad
        \text{error}_7=-0.0045,\quad
        \text{s.e.}_8=0.9587.
    \]
    \item Missing \(w[j]\) errors:
    \[
        w[0]\text{ error for ID 3}=0.0503,\quad
        w[2]\text{ error for ID 7}=-0.0001,\quad
        w[8]\text{ error for ID 2}=0.0000.
    \]
    \item LSE-updated weights:
    \[
    (0.3113,\;0.1835,\;-0.1448,\;-0.1588,\;0.1969,\;-0.6277,\;-0.7102,\;-0.4711,\;-0.1429).
    \]
    \item Using the LSE-updated weights, digits 7 and 8 both predict class \(0\):
    \[
        \sigma(z_7)=0.2680,\qquad \sigma(z_8)=0.0061.
    \]
    \item Using the MLE one-sample ID 0 update, the weights are effectively
    unchanged at four decimal places, and digits 7 and 8 again predict class \(0\):
    \[
        \sigma(z_7)=0.0139,\qquad \sigma(z_8)=0.0003.
    \]
\end{itemize}

\end{document}
